% Generated draft: requires mathematical and source review.
\documentclass[11pt]{article}
\usepackage[margin=25mm]{geometry}
\usepackage{fontspec}
\setmainfont{Latin Modern Roman}
\usepackage{amsmath,amssymb,mathtools,unicode-math}
\setmathfont{Latin Modern Math}
\usepackage{enumitem,needspace,xurl,hyperref}
\setlength{\parindent}{0pt}
\setlength{\parskip}{5pt}
\setlength{\emergencystretch}{3em}
\newcommand{\N}{\mathbb{N}}
\newcommand{\Z}{\mathbb{Z}}
\newcommand{\Q}{\mathbb{Q}}
\newcommand{\R}{\mathbb{R}}
\newcommand{\C}{\mathbb{C}}
\newcommand{\F}{\mathbb{F}}
\newcommand{\A}{\mathbb{A}}
\newcommand{\PP}{\mathbb{P}}
\newcommand{\E}{\mathbb{E}}
\newcommand{\eps}{\varepsilon}
\newcommand{\dd}{\,\mathrm{d}}
\begin{document}
\section*{UnsolvedMath ID: 30002659}
\textbf{Source problem code:} catalogue entry
\subsection{Definitions and mathematical statement}
Let \(d\geq 2\). A Euclidean convex body is a compact convex set \(K\subset\R^d\) with nonempty interior. For \(u\in\R^d\), its support function is
\[
h_K(u)=\max_{x\in K}\langle x,u\rangle.
\]
The body has constant width \(1\) if
\[
h_K(u)+h_K(-u)=1\qquad\text{for every }u\in\R^d\text{ with }\|u\|=1.
\]
No smoothness assumption is supplied.

To make the billiard terminology explicit, consider the following generalized reflection convention, whose agreement with the intended source formulation requires review. At \(x\in\partial K\), define the outward normal cone by
\[
N_K(x)=\{n\in\R^d:\langle n,y-x\rangle\leq 0\text{ for every }y\in K\}.
\]
A nonconstant closed generalized billiard trajectory is a cyclic polygon \(\gamma=(x_1,\ldots,x_m)\), where \(m\geq2\), each \(x_i\in\partial K\), consecutive vertices are distinct, and, with indices read modulo \(m\),
\[
v_i=\frac{x_{i+1}-x_i}{\|x_{i+1}-x_i\|},\qquad
v_{i-1}-v_i\in N_K(x_i)\setminus\{0\}.
\]
Its length is
\[
L(\gamma)=\sum_{i=1}^{m}\|x_{i+1}-x_i\|.
\]
It is shortest if \(L(\gamma)\leq L(\eta)\) for every nonconstant closed generalized billiard trajectory \(\eta\) in \(K\). Its period is the least number of reflections in one traversal of its primitive cyclic trajectory; period \(2\) means travel from one boundary point to another and back along the same segment.

\textbf{Question, using the universal interpretation in the supplied literature assessment.} For every integer \(d\geq3\), every Euclidean convex body \(K\subset\R^d\) of constant width \(1\), and every shortest nonconstant closed billiard trajectory \(\gamma\) in \(K\), must \(\gamma\) have period \(2\)?

The supplied record presents the planar period-two conclusion as known and asks for its extension to every higher dimension. The record does not specify a reflection convention at nonsmooth boundary points. Consequently, the generalized convention above is an explicit provisional interpretation, not an attributed definition from the report. There is also a quantifier issue: the statement says ``a shortest'' trajectory, whereas the supplied assessment explicitly asks whether \emph{every} shortest trajectory has period two. The weaker alternative would ask only whether each such body admits a shortest trajectory of period two; these targets must be distinguished in mathematical review.
\subsection{Short English statement}
The question asks whether, in every Euclidean dimension \(d\geq3\), every shortest closed billiard trajectory in every convex body of constant width \(1\) travels between two boundary points and back. The supplied record reports the planar result as known. The treatment of nonsmooth boundaries and the distinction between existence of a period-two minimizer and the period-two property for every minimizer require confirmation.
\subsection{Sources}
\begin{itemize}
\item Supplied local UnsolvedMath registry record, permanent ID \(30002659\), title \emph{Shortest Billiard Trajectories in Constant Width Bodies}, source problem code \texttt{catalogue entry}, in the Oberwolfach Reports Open Problems collection. The record labels the statement ``exact'' and says it was retained byte-for-byte; this is supplied registry metadata, not independent source verification. No individual proposer is supplied.
\item Supplied report citation: \emph{Discrete Geometry (2015)}, Oberwolfach Reports, DOI \texttt{10.4171/owr/2014/40}, pp.\ 2235--2296. Original Oberwolfach report: \url{https://doi.org/10.4171/owr/2014/40}. The record separately supplies proposed year \(2014\); the relationship between that year and the citation's \(2015\) label is not resolved here.
\item Supplied literature reference: Bezdek and Bezdek, shortest billiard trajectories, \url{https://doi.org/10.1007/s10711-009-9353-6}. The supplied record does not provide a more complete bibliographic citation or a theorem locator.
\item Literature assessment checked at \texttt{2026-08-21}: the supplied curation labels the problem open, reports the planar period-two result as known, and states that the assertion for every shortest trajectory in higher dimensions remains open. This is supplied dated curation derived from the OWR record, not independently verified literature assessment. The record explicitly requests independent verification of time-sensitive status and cited claims.
\end{itemize}
\end{document}
